% Drawing for measurement `jfet-tracer`
% Stands for the deck lines:
%   VK1 sw dk1 0
%   VGK1 gk1 0 {-k*step}
%   JK1 dk1 gk1 0 DUT
%   VSW sw 0 0
%   .dc VSW 0 25 0.02
\documentclass[border=8pt]{standalone}
\usepackage[siunitx,RPvoltages,european]{circuitikz}
\ctikzset{bipoles/length=0.9cm, label/align=straight, inductors/scale=0.8}
\ctikzset{inductor=american}
\usepackage{amsmath}
\begin{document}
\begin{circuitikz}[font=\footnotesize, line width=0.5pt]
\draw (0,0) node[njfet](Q){};
\node[anchor=west] at ($(Q.D)!0.5!(Q.S) + (0.55,0)$) {DUT};
\draw (Q.S) to[short] ++(0,-0.5) node[ground]{};
\draw (Q.D) to[short] ++(0,0.3) to[rmeter, t=A, l={$I_D$}] ++(0,1.4) to[short] ++(2.0,0) to[vsource, l={$V_{DS}$ $0\to\SI{25}{V}$}, invert] ++(0,-2.4) node[ground]{};
\node[anchor=west, align=left] at (2.3,2.1) {$\rightarrow$ swept};
\draw (Q.G) to[short] ++(-0.5,0) to[short] ++(-1.4,0) to[short] ++(-0.5,0) coordinate(gl) to[vsource, l_={\shortstack[r]{$V_{GS}=-k\cdot\text{step}$\\ $k=0\dots5$}}, invert] ++(0,-1.9) node[ground]{};
\node[anchor=west, align=left] at (-4.6,-3.6) {one copy per step; drawn as one\\ step: auto ($V_{GS(\mathrm{off})}/5$) or 0.05\dots\SI{3}{V}};
\end{circuitikz}
\end{document}
